\documentclass[12pt]{article}

% ============================================================
% IDIOMA Y CODIFICACIÓN
% ============================================================
\usepackage[utf8]{inputenc}
\usepackage[T1]{fontenc}
\usepackage[spanish, es-noshorthands]{babel}

% ============================================================
% COLOR
% ============================================================
\usepackage[dvipsnames]{xcolor}

% ============================================================
% MÁRGENES Y GEOMETRÍA
% ============================================================
\usepackage[a4paper, margin=2.5cm]{geometry}

% ============================================================
% ESPACIADO MEJORADO
% ============================================================
\usepackage{setspace}
\setstretch{1.15}

\setlength{\parskip}{0.5em}
\setlength{\parindent}{0pt}

% ============================================================
% FORMATO DE TÍTULOS
% ============================================================
\usepackage{titlesec}

\titleformat{\section}
  {\Large\bfseries}
  {\thesection.}
  {0.5em}
  {}

\titlespacing*{\section}{0pt}{1.5em}{0.8em}

% ============================================================
% MATEMÁTICAS
% ============================================================
\usepackage{amsmath, amssymb, amsfonts}

% ============================================================
% GRÁFICOS
% ============================================================
\usepackage{graphicx}
\usepackage{float}
\usepackage{tikz}
\usetikzlibrary{arrows.meta}
\usepackage{pgfplots}
\pgfplotsset{compat=1.18}

% ============================================================
% LISTAS Y NUMERACIÓN
% ============================================================
\usepackage{enumitem}

% ============================================================
% ENLACES E HIPERTEXTO
% ============================================================
\usepackage{hyperref}
\hypersetup{
    colorlinks  = true,
    linkcolor   = black,
    urlcolor    = blue!60!black,
    citecolor   = black,
    pdftitle    = {Trabajo Práctico N°4},
    pdfauthor   = {Facundo Luna},
}

% ============================================================
% ENCABEZADO Y PIE DE PÁGINA
% ============================================================
\usepackage{fancyhdr}
\pagestyle{fancy}
\fancyhf{}
\renewcommand{\headrulewidth}{0.4pt}
\renewcommand{\footrulewidth}{0.4pt}
\fancyhead[L]{\small Geometría Analítica}
\fancyhead[R]{\small Trabajo Práctico N°4}
\fancyfoot[L]{\small Facundo Luna}
\fancyfoot[R]{\small \thepage}

% ============================================================
% DATOS DEL TRABAJO
% ============================================================

\newcommand{\universidad}{Universidad Católica de Salta}
\newcommand{\carrera}{Licenciatura en Ciencia de Datos}
\newcommand{\materia}{Geometría Analítica}
\newcommand{\trabajo}{Trabajo Práctico N°4}
\newcommand{\profesor}{Lic. Mauro Vaca}
\newcommand{\alumno}{Facundo Luna}
\newcommand{\fecha}{\today}

\newcommand{\logoincluido}{%
    \includegraphics[width=3.5cm]{../../../assets/ucasal-logo.png}\\[0.5cm]%
}
\newcommand{\logoportada}{\logoincluido}

% ============================================================
% COMANDO SOLUCIÓN
% ============================================================
\newcommand{\solucion}{%
    \vspace{0.8em}%
    \textbf{Solución:}%
    \vspace{0.8em}%
}

% ============================================================
% COMANDO VOLVER AL ÍNDICE
% ============================================================
\newcommand{\volverindice}{%
    \vspace{1em}%
    \begin{flushright}%
        \hyperlink{indice}{\small $\uparrow$ Volver al índice}%
    \end{flushright}%
}

% ============================================================

\begin{document}

% ======================== PORTADA ===========================

\begin{titlepage}
    \centering

    \logoportada

    \rule{\linewidth}{1.5pt}\\[0.4cm]

    {\large\scshape \universidad}\\[0.2cm]
    {\large \carrera}\\[0.1cm]

    \rule{\linewidth}{0.5pt}\\[1.5cm]

    {\Huge\bfseries \trabajo}\\[0.5cm]
    {\LARGE \materia}\\[2cm]

    \begin{tabular}{rl}
        \textbf{Alumno:}   & \alumno   \\[0.3cm]
        \textbf{Profesor:} & \profesor \\[0.3cm]
        \textbf{Fecha:}    & \fecha    \\
    \end{tabular}

    \vfill

    \rule{\linewidth}{1.5pt}

\end{titlepage}

% ======================== ÍNDICE ============================

\newpage
\tableofcontents
\hypertarget{indice}{}
\setcounter{tocdepth}{2}


% ======================== EJERCICIO 1 ========================

\newpage

\section{Ecuación de la circunferencia}

Determine las ecuaciones de las circunferencias que satisfacen las siguientes
condiciones, luego grafíquelas.

\vspace{0.5em}
Usamos en todos los casos la \textbf{ecuación canónica} de la circunferencia de
centro $C(h,k)$ y radio $r$:
\[
(x-h)^{2}+(y-k)^{2}=r^{2}.
\]

\begin{enumerate}[label=\alph*.-]

    \item Su radio es $r=3$ y tiene centro en el origen de coordenadas.

    \solucion

    Con $C(0,0)$ y $r=3$:
    \[
    (x-0)^{2}+(y-0)^{2}=3^{2}.
    \]

    \[
    \boxed{
    x^{2}+y^{2}=9
    }
    \]

    En forma general: $x^{2}+y^{2}-9=0$.

    \begin{center}
    \begin{tikzpicture}
        \begin{axis}[
            width=8cm, axis equal image,
            axis lines=middle,
            xlabel={$x$}, ylabel={$y$},
            xmin=-4.5, xmax=4.5, ymin=-4.5, ymax=4.5,
            xtick={-3,-2,-1,1,2,3}, ytick={-3,-2,-1,1,2,3},
            axis line style={-Latex},
            every axis x label/.style={at={(ticklabel* cs:1.0)}, anchor=west},
            every axis y label/.style={at={(ticklabel* cs:1.0)}, anchor=south},
        ]
            \addplot[thick, blue, domain=0:360, samples=200]
                ({3*cos(x)}, {3*sin(x)});
            \addplot[only marks, mark=*, mark size=1.6pt, black]
                coordinates {(0,0)};
            \draw[dashed, red] (axis cs:0,0) -- (axis cs:2.1213,2.1213)
                node[midway, above left, inner sep=1pt] {$r=3$};
        \end{axis}
    \end{tikzpicture}
    \end{center}

    \item Tiene centro en el punto $(-1,2)$ y su diámetro es de $8$ unidades.

    \solucion

    Del diámetro obtenemos el radio:
    \[
    r=\frac{d}{2}=\frac82=4 .
    \]

    Con $C(-1,2)$ y $r=4$:
    \[
    \big(x-(-1)\big)^{2}+(y-2)^{2}=4^{2}.
    \]

    \[
    \boxed{
    (x+1)^{2}+(y-2)^{2}=16
    }
    \]

    Desarrollando se obtiene la forma general:
    \[
    x^{2}+2x+1+y^{2}-4y+4=16
    \quad\Longrightarrow\quad
    x^{2}+y^{2}+2x-4y-11=0 .
    \]

    \begin{center}
    \begin{tikzpicture}
        \begin{axis}[
            width=8.5cm, axis equal image,
            axis lines=middle,
            xlabel={$x$}, ylabel={$y$},
            xmin=-6, xmax=4, ymin=-3, ymax=7,
            xtick={-5,-4,-3,-2,-1,1,2,3}, ytick={-2,-1,1,2,3,4,5,6},
            axis line style={-Latex},
            every axis x label/.style={at={(ticklabel* cs:1.0)}, anchor=west},
            every axis y label/.style={at={(ticklabel* cs:1.0)}, anchor=south},
        ]
            \addplot[thick, blue, domain=0:360, samples=200]
                ({-1+4*cos(x)}, {2+4*sin(x)});
            \addplot[only marks, mark=*, mark size=1.6pt, black]
                coordinates {(-1,2)};
            \node[below left, inner sep=3pt] at (axis cs:-1,2) {$C$};
            \draw[dashed, red] (axis cs:-1,2) -- (axis cs:1.8284,4.8284)
                node[midway, above left, inner sep=1pt] {$r=4$};
        \end{axis}
    \end{tikzpicture}
    \end{center}

    \item El centro está en el punto $(-2,5)$ y pasa por el origen de
    coordenadas.

    \solucion

    El radio es la distancia del centro al punto por el que pasa la
    circunferencia:
    \[
    r=d\big(C,O\big)
    =\sqrt{\big(0-(-2)\big)^{2}+(0-5)^{2}}
    =\sqrt{4+25}
    =\sqrt{29}\approx 5{,}39 .
    \]

    Con $C(-2,5)$ y $r=\sqrt{29}$:

    \[
    \boxed{
    (x+2)^{2}+(y-5)^{2}=29
    }
    \]

    En forma general:
    \[
    x^{2}+4x+4+y^{2}-10y+25=29
    \quad\Longrightarrow\quad
    x^{2}+y^{2}+4x-10y=0 .
    \]

    El término independiente es nulo, lo que confirma que la curva pasa por el
    origen: al reemplazar $x=0$, $y=0$ la ecuación se satisface.

    \begin{center}
    \begin{tikzpicture}
        \begin{axis}[
            width=8.5cm, axis equal image,
            axis lines=middle,
            xlabel={$x$}, ylabel={$y$},
            xmin=-8.5, xmax=4.5, ymin=-1.5, ymax=11,
            xtick={-7,-5,-3,-1,1,3}, ytick={1,3,5,7,9},
            axis line style={-Latex},
            every axis x label/.style={at={(ticklabel* cs:1.0)}, anchor=west},
            every axis y label/.style={at={(ticklabel* cs:1.0)}, anchor=south},
        ]
            \addplot[thick, blue, domain=0:360, samples=200]
                ({-2+5.385164807*cos(x)}, {5+5.385164807*sin(x)});
            \addplot[only marks, mark=*, mark size=1.6pt, black]
                coordinates {(-2,5) (0,0)};
            \node[left, inner sep=5pt] at (axis cs:-2,5) {$C(-2,5)$};
            \node[above left, inner sep=3pt] at (axis cs:0,0) {$O$};
            \draw[dashed, red] (axis cs:-2,5) -- (axis cs:0,0)
                node[midway, left, inner sep=5pt] {$r=\sqrt{29}$};
        \end{axis}
    \end{tikzpicture}
    \end{center}

\end{enumerate}

\volverindice
\newpage

% ======================== EJERCICIO 2 ========================

\section{Elementos de la elipse a partir de su gráfica}

A partir de las siguientes gráficas, determine las coordenadas del centro,
semieje mayor y menor, ejes mayor y menor, vértices, distancia focal y ecuación
de la elipse.

\vspace{0.5em}
En cada caso leemos de la figura los cuatro puntos marcados: los dos
\emph{vértices} $V$ y $V'$ (extremos del eje mayor) y los dos extremos del eje
menor $M$ y $M'$. El centro es el punto medio de cualquiera de los dos ejes, y
la semidistancia focal se obtiene de la relación
\[
c^{2}=a^{2}-b^{2}.
\]

\begin{enumerate}[label=\alph*.-]

    \item Datos leídos de la gráfica:
    $V'(-7,1)$, $V(3,1)$, $M(-2,3)$, $M'(-2,-1)$.

    \solucion

    \textbf{1. Centro.} Es el punto medio del eje mayor $\overline{V'V}$:
    \[
    C=\left(\frac{-7+3}{2},\;\frac{1+1}{2}\right)=(-2,\,1).
    \]
    Se verifica con el eje menor:
    $\left(\frac{-2-2}{2},\frac{3-1}{2}\right)=(-2,1)$.

    \textbf{2. Semiejes.} Como $V'$ y $V$ tienen la misma ordenada, el eje mayor
    es \emph{horizontal}:
    \[
    2a=3-(-7)=10\;\Rightarrow\;a=5,
    \qquad
    2b=3-(-1)=4\;\Rightarrow\;b=2 .
    \]

    \textbf{3. Ejes.} El eje mayor mide $2a=10$ y está sobre la recta $y=1$; el
    eje menor mide $2b=4$ y está sobre la recta $x=-2$.

    \textbf{4. Vértices.} $V(3,1)$ y $V'(-7,1)$ sobre el eje mayor;
    $M(-2,3)$ y $M'(-2,-1)$ sobre el eje menor.

    \textbf{5. Distancia focal.}
    \[
    c^{2}=a^{2}-b^{2}=25-4=21
    \;\Rightarrow\;
    c=\sqrt{21}\approx 4{,}58 ,
    \]
    de modo que la distancia focal es
    \[
    2c=2\sqrt{21}\approx 9{,}17,
    \]
    y los focos son $F(-2+\sqrt{21},\,1)$ y $F'(-2-\sqrt{21},\,1)$.

    \textbf{6. Ecuación.} Con eje mayor horizontal:

    \[
    \boxed{
    \frac{(x+2)^{2}}{25}+\frac{(y-1)^{2}}{4}=1
    }
    \]

    La excentricidad es $e=\dfrac{c}{a}=\dfrac{\sqrt{21}}{5}\approx 0{,}92$.

    \begin{center}
    \begin{tikzpicture}
        \begin{axis}[
            width=12cm, axis equal image,
            axis lines=middle,
            xlabel={$x$}, ylabel={$y$},
            xmin=-9.5, xmax=5.5, ymin=-3.2, ymax=4.6,
            xtick={-7,-5,-3,-1,1,3}, ytick={-2,-1,1,2,3},
            axis line style={-Latex},
            every axis x label/.style={at={(ticklabel* cs:1.0)}, anchor=west},
            every axis y label/.style={at={(ticklabel* cs:1.0)}, anchor=south},
        ]
            \addplot[thick, blue, domain=0:360, samples=200]
                ({-2+5*cos(x)}, {1+2*sin(x)});
            \addplot[only marks, mark=*, mark size=1.6pt, black]
                coordinates {(-2,1) (3,1) (-7,1) (-2,3) (-2,-1)};
            \addplot[only marks, mark=*, mark size=1.6pt, red]
                coordinates {(2.583,1) (-6.583,1)};
            \node[above left, inner sep=2pt] at (axis cs:-2,1) {$C$};
            \node[right, inner sep=5pt] at (axis cs:3,1) {$V$};
            \node[left, inner sep=5pt]  at (axis cs:-7,1) {$V'$};
            \node[above, inner sep=4pt] at (axis cs:-2,3) {$M$};
            \node[below, inner sep=4pt] at (axis cs:-2,-1) {$M'$};
            \node[below, red, inner sep=5pt] at (axis cs:2.583,1) {$F$};
            \node[below, red, inner sep=5pt] at (axis cs:-6.583,1) {$F'$};
        \end{axis}
    \end{tikzpicture}
    \end{center}

    \item Datos leídos de la gráfica:
    $V(1,2)$, $V'(1,-6)$, $M(3,-2)$, $M'(-1,-2)$.

    \solucion

    \textbf{1. Centro.} Punto medio del eje mayor $\overline{VV'}$:
    \[
    C=\left(\frac{1+1}{2},\;\frac{2+(-6)}{2}\right)=(1,\,-2).
    \]
    Se verifica con el eje menor:
    $\left(\frac{3-1}{2},\frac{-2-2}{2}\right)=(1,-2)$.

    \textbf{2. Semiejes.} Como $V$ y $V'$ tienen la misma abscisa, el eje mayor
    es \emph{vertical}:
    \[
    2a=2-(-6)=8\;\Rightarrow\;a=4,
    \qquad
    2b=3-(-1)=4\;\Rightarrow\;b=2 .
    \]

    \textbf{3. Ejes.} El eje mayor mide $2a=8$ y está sobre la recta $x=1$; el
    eje menor mide $2b=4$ y está sobre la recta $y=-2$.

    \textbf{4. Vértices.} $V(1,2)$ y $V'(1,-6)$ sobre el eje mayor;
    $M(3,-2)$ y $M'(-1,-2)$ sobre el eje menor.

    \textbf{5. Distancia focal.}
    \[
    c^{2}=a^{2}-b^{2}=16-4=12
    \;\Rightarrow\;
    c=2\sqrt{3}\approx 3{,}46 ,
    \]
    de modo que la distancia focal es
    \[
    2c=4\sqrt{3}\approx 6{,}93,
    \]
    y los focos son $F(1,\,-2+2\sqrt{3})$ y $F'(1,\,-2-2\sqrt{3})$.

    \textbf{6. Ecuación.} Con eje mayor vertical, $a^{2}$ va debajo del término
    en $y$:

    \[
    \boxed{
    \frac{(x-1)^{2}}{4}+\frac{(y+2)^{2}}{16}=1
    }
    \]

    La excentricidad es
    $e=\dfrac{c}{a}=\dfrac{2\sqrt3}{4}=\dfrac{\sqrt3}{2}\approx 0{,}87$.

    \begin{center}
    \begin{tikzpicture}
        \begin{axis}[
            width=8cm, axis equal image,
            axis lines=middle,
            xlabel={$x$}, ylabel={$y$},
            xmin=-4.5, xmax=6.5, ymin=-7.8, ymax=3.8,
            xtick={-3,-1,1,3,5}, ytick={-7,-5,-3,-1,1,2},
            axis line style={-Latex},
            every axis x label/.style={at={(ticklabel* cs:1.0)}, anchor=west},
            every axis y label/.style={at={(ticklabel* cs:1.0)}, anchor=south},
        ]
            \addplot[thick, blue, domain=0:360, samples=200]
                ({1+2*cos(x)}, {-2+4*sin(x)});
            \addplot[only marks, mark=*, mark size=1.6pt, black]
                coordinates {(1,-2) (1,2) (1,-6) (3,-2) (-1,-2)};
            \addplot[only marks, mark=*, mark size=1.6pt, red]
                coordinates {(1,1.464) (1,-5.464)};
            \node[below left, inner sep=2pt] at (axis cs:1,-2) {$C$};
            \node[above, inner sep=4pt] at (axis cs:1,2) {$V$};
            \node[below, inner sep=4pt] at (axis cs:1,-6) {$V'$};
            \node[right, inner sep=5pt] at (axis cs:3,-2) {$M$};
            \node[left, inner sep=5pt]  at (axis cs:-1,-2) {$M'$};
            \node[right, red, inner sep=5pt] at (axis cs:1,1.464) {$F$};
            \node[right, red, inner sep=5pt] at (axis cs:1,-5.464) {$F'$};
        \end{axis}
    \end{tikzpicture}
    \end{center}

\end{enumerate}

\volverindice
\newpage

% ======================== EJERCICIO 3 ========================

\section{Elementos de la elipse a partir de su ecuación}

Determine todos los elementos y grafique las siguientes elipses.

\vspace{0.5em}
El procedimiento es siempre el mismo: agrupar los términos en $x$ y en $y$,
\emph{completar cuadrados} y dividir para llevar la ecuación a la forma
canónica
\[
\frac{(x-h)^{2}}{a^{2}}+\frac{(y-k)^{2}}{b^{2}}=1 ,
\]
donde el \emph{mayor} de los dos denominadores corresponde al eje mayor.

\begin{enumerate}[label=\alph*.-]

    \setcounter{enumi}{2} % comienza en c)

    \item $9x^{2}+8y^{2}-54x-16y+17=0$

    \solucion

    \textbf{1. Completamos cuadrados.} Agrupamos y extraemos los coeficientes
    de los términos cuadráticos:
    \[
    9\left(x^{2}-6x\right)+8\left(y^{2}-2y\right)+17=0 .
    \]

    Completando cada paréntesis:
    \[
    9\left[(x-3)^{2}-9\right]+8\left[(y-1)^{2}-1\right]+17=0 ,
    \]
    \[
    9(x-3)^{2}-81+8(y-1)^{2}-8+17=0 ,
    \]
    \[
    9(x-3)^{2}+8(y-1)^{2}=72 .
    \]

    \textbf{2. Forma canónica.} Dividimos por $72$:
    \[
    \boxed{
    \frac{(x-3)^{2}}{8}+\frac{(y-1)^{2}}{9}=1
    }
    \]

    \textbf{3. Elementos.} El mayor denominador ($9$) está bajo el término en
    $y$, de modo que el \textbf{eje mayor es vertical}:
    \[
    a^{2}=9\Rightarrow a=3,
    \qquad
    b^{2}=8\Rightarrow b=2\sqrt2\approx 2{,}83,
    \]
    \[
    c^{2}=a^{2}-b^{2}=9-8=1\Rightarrow c=1 .
    \]

    \begin{itemize}[leftmargin=2em, itemsep=0.1em]
        \item \textbf{Centro:} $C(3,1)$.
        \item \textbf{Eje mayor:} vertical, sobre la recta $x=3$; mide
        $2a=6$.
        \item \textbf{Eje menor:} horizontal, sobre la recta $y=1$; mide
        $2b=4\sqrt2\approx 5{,}66$.
        \item \textbf{Vértices:} $V(3,4)$ y $V'(3,-2)$.
        \item \textbf{Extremos del eje menor:}
        $M(3+2\sqrt2,\,1)\approx(5{,}83;\,1)$ y
        $M'(3-2\sqrt2,\,1)\approx(0{,}17;\,1)$.
        \item \textbf{Focos:} $F(3,2)$ y $F'(3,0)$; distancia focal $2c=2$.
        \item \textbf{Excentricidad:} $e=\dfrac{c}{a}=\dfrac13\approx 0{,}33$.
        \item \textbf{Lado recto:} $\dfrac{2b^{2}}{a}=\dfrac{16}{3}
        \approx 5{,}33$.
    \end{itemize}

    \begin{center}
    \begin{tikzpicture}
        \begin{axis}[
            width=8cm, axis equal image,
            axis lines=middle,
            xlabel={$x$}, ylabel={$y$},
            xmin=-1.5, xmax=7.5, ymin=-3.5, ymax=5.5,
            xtick={-1,1,2,3,4,5,6,7}, ytick={-3,-2,-1,1,2,3,4,5},
            axis line style={-Latex},
            every axis x label/.style={at={(ticklabel* cs:1.0)}, anchor=west},
            every axis y label/.style={at={(ticklabel* cs:1.0)}, anchor=south},
        ]
            \addplot[thick, blue, domain=0:360, samples=200]
                ({3+2.828427*cos(x)}, {1+3*sin(x)});
            \addplot[only marks, mark=*, mark size=1.6pt, black]
                coordinates {(3,1) (3,4) (3,-2) (5.828,1) (0.172,1)};
            \addplot[only marks, mark=*, mark size=1.6pt, red]
                coordinates {(3,2) (3,0)};
            \node[left, inner sep=5pt]  at (axis cs:3,1) {$C$};
            \node[above, inner sep=4pt] at (axis cs:3,4) {$V$};
            \node[below, inner sep=4pt] at (axis cs:3,-2) {$V'$};
            \node[right, inner sep=5pt] at (axis cs:5.828,1) {$M$};
            \node[left, inner sep=5pt]  at (axis cs:0.172,1) {$M'$};
            \node[right, red, inner sep=5pt] at (axis cs:3,2) {$F$};
            \node[right, red, inner sep=5pt] at (axis cs:3,0) {$F'$};
        \end{axis}
    \end{tikzpicture}
    \end{center}

    \item $4x^{2}+9y^{2}-8x-54y+49=0$

    \solucion

    \textbf{1. Completamos cuadrados.}
    \[
    4\left(x^{2}-2x\right)+9\left(y^{2}-6y\right)+49=0 ,
    \]
    \[
    4\left[(x-1)^{2}-1\right]+9\left[(y-3)^{2}-9\right]+49=0 ,
    \]
    \[
    4(x-1)^{2}-4+9(y-3)^{2}-81+49=0 ,
    \]
    \[
    4(x-1)^{2}+9(y-3)^{2}=36 .
    \]

    \textbf{2. Forma canónica.} Dividimos por $36$:
    \[
    \boxed{
    \frac{(x-1)^{2}}{9}+\frac{(y-3)^{2}}{4}=1
    }
    \]

    \textbf{3. Elementos.} Ahora el mayor denominador ($9$) está bajo el
    término en $x$, de modo que el \textbf{eje mayor es horizontal}:
    \[
    a^{2}=9\Rightarrow a=3,
    \qquad
    b^{2}=4\Rightarrow b=2,
    \qquad
    c^{2}=9-4=5\Rightarrow c=\sqrt5\approx 2{,}24 .
    \]

    \begin{itemize}[leftmargin=2em, itemsep=0.1em]
        \item \textbf{Centro:} $C(1,3)$.
        \item \textbf{Eje mayor:} horizontal, sobre la recta $y=3$; mide
        $2a=6$.
        \item \textbf{Eje menor:} vertical, sobre la recta $x=1$; mide $2b=4$.
        \item \textbf{Vértices:} $V(4,3)$ y $V'(-2,3)$.
        \item \textbf{Extremos del eje menor:} $M(1,5)$ y $M'(1,1)$.
        \item \textbf{Focos:} $F(1+\sqrt5,\,3)\approx(3{,}24;\,3)$ y
        $F'(1-\sqrt5,\,3)\approx(-1{,}24;\,3)$; distancia focal
        $2c=2\sqrt5\approx 4{,}47$.
        \item \textbf{Excentricidad:}
        $e=\dfrac{\sqrt5}{3}\approx 0{,}75$.
        \item \textbf{Lado recto:} $\dfrac{2b^{2}}{a}=\dfrac83\approx 2{,}67$.
    \end{itemize}

    \begin{center}
    \begin{tikzpicture}
        \begin{axis}[
            width=9cm, axis equal image,
            axis lines=middle,
            xlabel={$x$}, ylabel={$y$},
            xmin=-4.5, xmax=6.5, ymin=-1, ymax=6.5,
            xtick={-4,-3,-2,-1,1,2,3,4,5,6}, ytick={1,2,3,4,5,6},
            axis line style={-Latex},
            every axis x label/.style={at={(ticklabel* cs:1.0)}, anchor=west},
            every axis y label/.style={at={(ticklabel* cs:1.0)}, anchor=south},
        ]
            \addplot[thick, blue, domain=0:360, samples=200]
                ({1+3*cos(x)}, {3+2*sin(x)});
            \addplot[only marks, mark=*, mark size=1.6pt, black]
                coordinates {(1,3) (4,3) (-2,3) (1,5) (1,1)};
            \addplot[only marks, mark=*, mark size=1.6pt, red]
                coordinates {(3.236,3) (-1.236,3)};
            \node[below left, inner sep=2pt] at (axis cs:1,3) {$C$};
            \node[right, inner sep=5pt] at (axis cs:4,3) {$V$};
            \node[left, inner sep=5pt]  at (axis cs:-2,3) {$V'$};
            \node[above, inner sep=4pt] at (axis cs:1,5) {$M$};
            \node[below, inner sep=4pt] at (axis cs:1,1) {$M'$};
            \node[above, red, inner sep=4pt] at (axis cs:3.236,3) {$F$};
            \node[above, red, inner sep=4pt] at (axis cs:-1.236,3) {$F'$};
        \end{axis}
    \end{tikzpicture}
    \end{center}

\end{enumerate}

\volverindice
\newpage

% ======================== EJERCICIO 4 ========================

\section{Elementos de la parábola a partir de su ecuación}

Determine todos los elementos de las siguientes parábolas, luego grafíquelas.

\vspace{0.5em}
Llevamos cada ecuación a una de las dos formas canónicas
\[
(x-h)^{2}=4p\,(y-k)
\qquad\text{(eje de simetría vertical)},
\]
\[
(y-k)^{2}=4p\,(x-h)
\qquad\text{(eje de simetría horizontal)},
\]
donde $V(h,k)$ es el vértice y $|p|$ la distancia del vértice al foco y a la
directriz. El signo de $p$ indica hacia dónde abre la parábola, y la
\emph{longitud del lado recto} es $|4p|$.

\begin{enumerate}[label=\alph*.-]

    \setcounter{enumi}{2} % comienza en c)

    \item $(x+2)^{2}=\dfrac12(y-1)$

    \solucion

    La ecuación ya está en forma canónica, con eje de simetría vertical:
    \[
    \big(x-(-2)\big)^{2}=4p\,(y-1),
    \qquad
    4p=\frac12
    \;\Longrightarrow\;
    p=\frac18 .
    \]

    Como $p>0$, la parábola \textbf{abre hacia arriba}.

    \begin{itemize}[leftmargin=2em, itemsep=0.1em]
        \item \textbf{Vértice:} $V(-2,1)$.
        \item \textbf{Eje de simetría:} la recta vertical $x=-2$.
        \item \textbf{Foco:} $F\left(-2,\;1+\dfrac18\right)
        =\left(-2,\;\dfrac98\right)$.
        \item \textbf{Directriz:} la recta horizontal
        $y=1-\dfrac18=\dfrac78$.
        \item \textbf{Lado recto:} $|4p|=\dfrac12$; sus extremos son
        $\left(-2\pm\dfrac14,\;\dfrac98\right)$, es decir
        $\left(-\dfrac74,\dfrac98\right)$ y
        $\left(-\dfrac94,\dfrac98\right)$.
        \item \textbf{Concavidad:} hacia arriba.
    \end{itemize}

    Despejando, $y=1+2(x+2)^{2}$: la parábola es muy cerrada, por lo que en la
    gráfica el foco y la directriz quedan muy próximos al vértice.

    \begin{center}
    \begin{tikzpicture}
        \begin{axis}[
            width=9cm, height=7.5cm,
            axis lines=middle,
            xlabel={$x$}, ylabel={$y$},
            xmin=-4.5, xmax=0.5, ymin=-1, ymax=9.5,
            xtick={-4,-3,-2,-1}, ytick={1,2,4,6,8},
            axis line style={-Latex},
            every axis x label/.style={at={(ticklabel* cs:1.0)}, anchor=west},
            every axis y label/.style={at={(ticklabel* cs:1.0)}, anchor=south},
        ]
            \addplot[thick, blue, domain=-4:0, samples=120]
                {1+2*(x+2)^2};
            \addplot[dashed, OliveGreen, domain=-4.5:0.5] {0.875};
            \addplot[only marks, mark=*, mark size=1.6pt, black]
                coordinates {(-2,1)};
            \node[left, inner sep=5pt] at (axis cs:-2,1) {$V$};
            \draw[dashed, gray] (axis cs:-2,-1) -- (axis cs:-2,9.5);
        \end{axis}
    \end{tikzpicture}
    \end{center}

    Como el foco y la directriz distan sólo $\tfrac18$ del vértice, no se
    distinguen en la vista anterior. El siguiente detalle amplía el entorno del
    vértice:

    \begin{center}
    \begin{tikzpicture}
        \begin{axis}[
            width=9cm, height=6.5cm,
            axis lines=box,
            grid=both, grid style={gray!20},
            xlabel={$x$}, ylabel={$y$},
            xmin=-2.7, xmax=-1.3, ymin=0.78, ymax=1.5,
            xtick={-2.5,-2.25,-2,-1.75,-1.5},
            ytick={0.875,1,1.125,1.25,1.375},
            yticklabels={$\frac78$,$1$,$\frac98$,$\frac54$,$\frac{11}8$},
            tick label style={font=\small},
        ]
            \addplot[thick, blue, domain=-2.7:-1.3, samples=120]
                {1+2*(x+2)^2};
            \addplot[dashed, OliveGreen, domain=-2.7:-1.3] {0.875};
            \addplot[dashed, red] coordinates {(-2.25,1.125) (-1.75,1.125)};
            \addplot[only marks, mark=*, mark size=1.6pt, black]
                coordinates {(-2,1) (-2.25,1.125) (-1.75,1.125)};
            \addplot[only marks, mark=*, mark size=1.8pt, red]
                coordinates {(-2,1.125)};
            \node[below left, inner sep=2pt] at (axis cs:-2,1) {$V$};
            \node[above, red, inner sep=3pt] at (axis cs:-2,1.125) {$F$};
            \node[above right, OliveGreen, inner sep=2pt]
                at (axis cs:-2.65,0.885) {$y=\tfrac78$ (directriz)};
            \node[below right, red, inner sep=2pt]
                at (axis cs:-2.25,1.125) {lado recto};
            \draw[dashed, gray] (axis cs:-2,0.78) -- (axis cs:-2,1.5);
        \end{axis}
    \end{tikzpicture}
    \end{center}

    \setcounter{enumi}{4} % salta a e)

    \item $y^{2}-6y-12x-15=0$

    \solucion

    \textbf{1. Completamos cuadrados en $y$.}
    \[
    y^{2}-6y=12x+15 ,
    \]
    \[
    (y-3)^{2}-9=12x+15 ,
    \]
    \[
    (y-3)^{2}=12x+24=12\,(x+2).
    \]

    \textbf{2. Identificación.} La forma es
    $(y-k)^{2}=4p\,(x-h)$ con
    \[
    h=-2,\qquad k=3,\qquad 4p=12\;\Longrightarrow\;p=3 .
    \]

    Como $p>0$, la parábola \textbf{abre hacia la derecha}.

    \[
    \boxed{
    (y-3)^{2}=12\,(x+2)
    }
    \]

    \begin{itemize}[leftmargin=2em, itemsep=0.1em]
        \item \textbf{Vértice:} $V(-2,3)$.
        \item \textbf{Eje de simetría:} la recta horizontal $y=3$.
        \item \textbf{Foco:} $F(-2+3,\;3)=(1,3)$.
        \item \textbf{Directriz:} la recta vertical $x=-2-3=-5$.
        \item \textbf{Lado recto:} $|4p|=12$; sus extremos son $(1,9)$ y
        $(1,-3)$.
        \item \textbf{Concavidad:} hacia la derecha.
    \end{itemize}

    \begin{center}
    \begin{tikzpicture}
        \begin{axis}[
            width=9.5cm, axis equal image,
            axis lines=middle,
            xlabel={$x$}, ylabel={$y$},
            xmin=-6.5, xmax=6.5, ymin=-5.5, ymax=11.5,
            xtick={-5,-3,-1,1,3,5}, ytick={-3,3,6,9},
            axis line style={-Latex},
            every axis x label/.style={at={(ticklabel* cs:1.0)}, anchor=west},
            every axis y label/.style={at={(ticklabel* cs:1.0)}, anchor=south},
        ]
            \addplot[thick, blue, domain=-8.5:8.5, samples=150]
                ({-2+x^2/12}, {3+x});
            \addplot[dashed, OliveGreen] coordinates {(-5,-5.5) (-5,11.5)};
            \addplot[dashed, gray] coordinates {(-6.5,3) (6.5,3)};
            \addplot[only marks, mark=*, mark size=1.6pt, black]
                coordinates {(-2,3) (1,9) (1,-3)};
            \addplot[only marks, mark=*, mark size=1.6pt, red]
                coordinates {(1,3)};
            \node[above left, inner sep=2pt] at (axis cs:-2,3) {$V$};
            \node[above right, red, inner sep=2pt] at (axis cs:1,3) {$F$};
            \node[right, OliveGreen, inner sep=4pt]
                at (axis cs:-5,-4.8) {$x=-5$};
        \end{axis}
    \end{tikzpicture}
    \end{center}

    \item $2y^{2}-8x-24y+56=0$

    \solucion

    \textbf{1. Dividimos por $2$ y completamos cuadrados.}
    \[
    y^{2}-4x-12y+28=0 ,
    \]
    \[
    y^{2}-12y=4x-28 ,
    \]
    \[
    (y-6)^{2}-36=4x-28 ,
    \]
    \[
    (y-6)^{2}=4x+8=4\,(x+2).
    \]

    \textbf{2. Identificación.} Con $(y-k)^{2}=4p\,(x-h)$:
    \[
    h=-2,\qquad k=6,\qquad 4p=4\;\Longrightarrow\;p=1 .
    \]

    Como $p>0$, la parábola \textbf{abre hacia la derecha}.

    \[
    \boxed{
    (y-6)^{2}=4\,(x+2)
    }
    \]

    \begin{itemize}[leftmargin=2em, itemsep=0.1em]
        \item \textbf{Vértice:} $V(-2,6)$.
        \item \textbf{Eje de simetría:} la recta horizontal $y=6$.
        \item \textbf{Foco:} $F(-2+1,\;6)=(-1,6)$.
        \item \textbf{Directriz:} la recta vertical $x=-2-1=-3$.
        \item \textbf{Lado recto:} $|4p|=4$; sus extremos son $(-1,8)$ y
        $(-1,4)$.
        \item \textbf{Concavidad:} hacia la derecha.
    \end{itemize}

    \begin{center}
    \begin{tikzpicture}
        \begin{axis}[
            width=9.5cm, axis equal image,
            axis lines=middle,
            xlabel={$x$}, ylabel={$y$},
            xmin=-4.5, xmax=8.5, ymin=-0.5, ymax=13,
            xtick={-3,-1,1,3,5,7}, ytick={2,4,6,8,10,12},
            axis line style={-Latex},
            every axis x label/.style={at={(ticklabel* cs:1.0)}, anchor=west},
            every axis y label/.style={at={(ticklabel* cs:1.0)}, anchor=south},
        ]
            \addplot[thick, blue, domain=-6:6, samples=150]
                ({-2+x^2/4}, {6+x});
            \addplot[dashed, OliveGreen] coordinates {(-3,-0.5) (-3,13)};
            \addplot[dashed, gray] coordinates {(-4.5,6) (8.5,6)};
            \addplot[only marks, mark=*, mark size=1.6pt, black]
                coordinates {(-2,6) (-1,8) (-1,4)};
            \addplot[only marks, mark=*, mark size=1.6pt, red]
                coordinates {(-1,6)};
            \node[above left, inner sep=2pt] at (axis cs:-2,6) {$V$};
            \node[below right, red, inner sep=2pt] at (axis cs:-1,6) {$F$};
            \node[right, OliveGreen, inner sep=4pt]
                at (axis cs:-3,1.2) {$x=-3$};
        \end{axis}
    \end{tikzpicture}
    \end{center}

\end{enumerate}

\volverindice
\newpage

% ======================== EJERCICIO 5 ========================

\section{Ecuación de la parábola a partir de sus condiciones}

Halle la ecuación de la parábola que cumple con las condiciones dadas, luego
grafíquela.

\begin{enumerate}[label=\alph*.-]

    \setcounter{enumi}{1} % comienza en b)

    \item Tiene foco en el punto $F=(0,-4)$ y directriz la recta $y-4=0$.

    \solucion

    \textbf{1. Orientación.} La directriz es la recta horizontal $y=4$, de modo
    que el eje de simetría es perpendicular a ella y pasa por el foco: es la
    recta vertical $x=0$. La ecuación tiene entonces la forma
    \[
    (x-h)^{2}=4p\,(y-k).
    \]

    \textbf{2. Vértice.} El vértice es el punto medio entre el foco y el pie de
    la directriz sobre el eje de simetría, es decir entre $(0,-4)$ y $(0,4)$:
    \[
    V=\left(\frac{0+0}{2},\;\frac{-4+4}{2}\right)=(0,0).
    \]

    \textbf{3. Parámetro.} $p$ es la distancia orientada del vértice al foco:
    \[
    p=-4-0=-4
    \quad\Longrightarrow\quad
    4p=-16 .
    \]
    El signo negativo indica que la parábola \textbf{abre hacia abajo}, como era
    de esperar porque el foco está por debajo de la directriz.

    \textbf{4. Ecuación.}

    \[
    \boxed{
    x^{2}=-16y
    }
    \]

    \textbf{Verificación:} con $4p=-16$ la directriz es $y=k-p=0-(-4)=4$
    \checkmark, y el foco es $(0,\,k+p)=(0,-4)$ \checkmark.

    \textbf{Elementos:} $V(0,0)$; $F(0,-4)$; directriz $y=4$; eje de simetría
    $x=0$; lado recto $|4p|=16$, con extremos $(8,-4)$ y $(-8,-4)$.

    \begin{center}
    \begin{tikzpicture}
        \begin{axis}[
            width=10cm, axis equal image,
            axis lines=middle,
            xlabel={$x$}, ylabel={$y$},
            xmin=-11, xmax=11, ymin=-8.5, ymax=6,
            xtick={-8,-4,4,8}, ytick={-8,-4,4},
            axis line style={-Latex},
            every axis x label/.style={at={(ticklabel* cs:1.0)}, anchor=west},
            every axis y label/.style={at={(ticklabel* cs:1.0)}, anchor=south},
        ]
            \addplot[thick, blue, domain=-10.5:10.5, samples=150]
                {-x^2/16};
            \addplot[dashed, OliveGreen] coordinates {(-11,4) (11,4)};
            \addplot[dashed, gray] coordinates {(0,-8.5) (0,6)};
            \addplot[only marks, mark=*, mark size=1.6pt, black]
                coordinates {(0,0) (8,-4) (-8,-4)};
            \addplot[only marks, mark=*, mark size=1.6pt, red]
                coordinates {(0,-4)};
            \node[above right] at (axis cs:0,0) {$V(0,0)$};
            \node[right, red]  at (axis cs:0,-4) {$F(0,-4)$};
            \node[above right, OliveGreen] at (axis cs:-10.5,4) {$y=4$};
        \end{axis}
    \end{tikzpicture}
    \end{center}

    \item Tiene vértice en el origen de coordenadas, es simétrica al eje de las
    ordenadas y pasa por el punto $(2,-3)$.

    \solucion

    \textbf{1. Forma de la ecuación.} Ser simétrica respecto del eje de las
    ordenadas (el eje $y$) significa que ese eje es el eje de simetría. Con
    vértice en el origen, la ecuación es
    \[
    x^{2}=4p\,y .
    \]

    \textbf{2. Determinación de $p$.} Imponemos que el punto $(2,-3)$ pertenezca
    a la parábola:
    \[
    2^{2}=4p\,(-3)
    \quad\Longrightarrow\quad
    4=-12p
    \quad\Longrightarrow\quad
    p=-\frac13 ,
    \qquad
    4p=-\frac43 .
    \]

    Como $p<0$, la parábola \textbf{abre hacia abajo}, lo que es coherente con
    que pase por un punto de ordenada negativa.

    \textbf{3. Ecuación.}

    \[
    \boxed{
    x^{2}=-\frac43\,y
    \qquad\text{equivalentemente}\qquad
    3x^{2}+4y=0
    }
    \]

    \textbf{Verificación:} $3(2)^{2}+4(-3)=12-12=0$ \checkmark.

    \textbf{Elementos:} $V(0,0)$; $F\left(0,-\dfrac13\right)$; directriz
    $y=\dfrac13$; eje de simetría $x=0$; lado recto $|4p|=\dfrac43$.

    \begin{center}
    \begin{tikzpicture}
        \begin{axis}[
            width=9cm, axis equal image,
            axis lines=middle,
            xlabel={$x$}, ylabel={$y$},
            xmin=-4.2, xmax=5.6, ymin=-7.5, ymax=2,
            xtick={-4,-3,-2,-1,1,2,3,4,5}, ytick={-6,-4,-3,-2,1},
            axis line style={-Latex},
            every axis x label/.style={at={(ticklabel* cs:1.0)}, anchor=west},
            every axis y label/.style={at={(ticklabel* cs:1.0)}, anchor=south},
        ]
            \addplot[thick, blue, domain=-3.1:3.1, samples=150]
                {-3*x^2/4};
            \addplot[dashed, OliveGreen] coordinates {(-3.8,0.3333) (3.8,0.3333)};
            \addplot[dashed, gray] coordinates {(0,-7.5) (0,2)};
            \addplot[only marks, mark=*, mark size=1.6pt, black]
                coordinates {(0,0) (2,-3)};
            \addplot[only marks, mark=*, mark size=1.6pt, red]
                coordinates {(0,-0.3333)};
            \node[below left, inner sep=3pt] at (axis cs:0,0) {$V$};
            \draw[red, thin] (axis cs:0.05,-0.40) -- (axis cs:0.55,-1.15);
            \node[right, red, inner sep=1pt]
                at (axis cs:0.5,-1.25) {$F\left(0,-\tfrac13\right)$};
            \node[right, inner sep=5pt] at (axis cs:2,-3) {$P(2,-3)$};
            \node[above right, OliveGreen, inner sep=2pt]
                at (axis cs:-4.15,0.3333) {$y=\tfrac13$};
        \end{axis}
    \end{tikzpicture}
    \end{center}

\end{enumerate}

\volverindice
\newpage

% ======================== EJERCICIO 6 ========================

\section{Gráfica y elementos de la hipérbola}

Grafique las siguientes hipérbolas, señalando sus elementos.

\vspace{0.5em}
Llevamos cada ecuación a la forma canónica. Si el término positivo es el de
$x$,
\[
\frac{(x-h)^{2}}{a^{2}}-\frac{(y-k)^{2}}{b^{2}}=1 ,
\]
el eje transverso es horizontal; si el positivo es el de $y$, es vertical. En
ambos casos
\[
c^{2}=a^{2}+b^{2},
\]
y las asíntotas pasan por el centro con pendientes $\pm\dfrac{b}{a}$
(eje transverso horizontal).

\begin{enumerate}[label=\alph*.-]

    \setcounter{enumi}{2} % comienza en c)

    \item $4x^{2}-y^{2}-16x+4y+8=0$

    \solucion

    \textbf{1. Completamos cuadrados.}
    \[
    4\left(x^{2}-4x\right)-\left(y^{2}-4y\right)+8=0 ,
    \]
    \[
    4\left[(x-2)^{2}-4\right]-\left[(y-2)^{2}-4\right]+8=0 ,
    \]
    \[
    4(x-2)^{2}-16-(y-2)^{2}+4+8=0 ,
    \]
    \[
    4(x-2)^{2}-(y-2)^{2}=4 .
    \]

    \textbf{2. Forma canónica.} Dividimos por $4$:
    \[
    \boxed{
    \frac{(x-2)^{2}}{1}-\frac{(y-2)^{2}}{4}=1
    }
    \]

    \textbf{3. Elementos.} El término positivo es el de $x$, de modo que el
    \textbf{eje transverso es horizontal}:
    \[
    a^{2}=1\Rightarrow a=1,
    \qquad
    b^{2}=4\Rightarrow b=2,
    \qquad
    c^{2}=1+4=5\Rightarrow c=\sqrt5\approx 2{,}24 .
    \]

    \begin{itemize}[leftmargin=2em, itemsep=0.1em]
        \item \textbf{Centro:} $C(2,2)$.
        \item \textbf{Vértices:} $V(3,2)$ y $V'(1,2)$; eje transverso $2a=2$.
        \item \textbf{Eje conjugado:} $2b=4$, sobre la recta $x=2$.
        \item \textbf{Focos:} $F(2+\sqrt5,\,2)\approx(4{,}24;\,2)$ y
        $F'(2-\sqrt5,\,2)\approx(-0{,}24;\,2)$; distancia focal
        $2c=2\sqrt5\approx 4{,}47$.
        \item \textbf{Asíntotas:} $y-2=\pm\dfrac{b}{a}(x-2)=\pm2(x-2)$, es
        decir
        \[
        y=2x-2
        \qquad\text{y}\qquad
        y=-2x+6 .
        \]
        \item \textbf{Excentricidad:} $e=\dfrac{c}{a}=\sqrt5\approx 2{,}24$.
        \item \textbf{Lado recto:} $\dfrac{2b^{2}}{a}=8$.
    \end{itemize}

    \begin{center}
    \begin{tikzpicture}
        \begin{axis}[
            width=8.5cm, axis equal image,
            axis lines=middle,
            xlabel={$x$}, ylabel={$y$},
            xmin=-2.8, xmax=6.8, ymin=-3.6, ymax=7.2,
            xtick={-2,-1,1,2,3,4,5,6}, ytick={-3,-2,-1,2,4,6},
            axis line style={-Latex},
            every axis x label/.style={at={(ticklabel* cs:1.0)}, anchor=west},
            every axis y label/.style={at={(ticklabel* cs:1.0)}, anchor=south},
        ]
            \addplot[dashed, OliveGreen, domain=-3.5:7.5] {2*x-2};
            \addplot[dashed, OliveGreen, domain=-3.5:7.5] {-2*x+6};
            \addplot[thick, blue, domain=-1.8:1.8, samples=120]
                ({2+cosh(x)}, {2+2*sinh(x)});
            \addplot[thick, blue, domain=-1.8:1.8, samples=120]
                ({2-cosh(x)}, {2+2*sinh(x)});
            \addplot[only marks, mark=*, mark size=1.6pt, black]
                coordinates {(2,2) (3,2) (1,2)};
            \addplot[only marks, mark=*, mark size=1.6pt, red]
                coordinates {(4.236,2) (-0.236,2)};
            \node[below, inner sep=4pt] at (axis cs:2,2) {$C$};
            \node[above right, inner sep=1pt] at (axis cs:3,2) {$V$};
            \node[above left, inner sep=1pt]  at (axis cs:1,2) {$V'$};
            \node[above right, red, inner sep=1pt] at (axis cs:4.236,2) {$F$};
            \node[above left, red, inner sep=1pt] at (axis cs:-0.236,2) {$F'$};
        \end{axis}
    \end{tikzpicture}
    \end{center}

    \item $-x^{2}+16y^{2}-8x-128y+256=0$

    \solucion

    \textbf{1. Completamos cuadrados.}
    \[
    -\left(x^{2}+8x\right)+16\left(y^{2}-8y\right)+256=0 ,
    \]
    \[
    -\left[(x+4)^{2}-16\right]+16\left[(y-4)^{2}-16\right]+256=0 ,
    \]
    \[
    -(x+4)^{2}+16+16(y-4)^{2}-256+256=0 ,
    \]
    \[
    16(y-4)^{2}-(x+4)^{2}=-16 .
    \]

    \textbf{2. Forma canónica.} Multiplicamos por $-1$ y dividimos por $16$:
    \[
    (x+4)^{2}-16(y-4)^{2}=16 ,
    \]
    \[
    \boxed{
    \frac{(x+4)^{2}}{16}-\frac{(y-4)^{2}}{1}=1
    }
    \]

    \textbf{3. Elementos.} El término positivo es el de $x$, de modo que el
    \textbf{eje transverso es horizontal}:
    \[
    a^{2}=16\Rightarrow a=4,
    \qquad
    b^{2}=1\Rightarrow b=1,
    \qquad
    c^{2}=16+1=17\Rightarrow c=\sqrt{17}\approx 4{,}12 .
    \]

    \begin{itemize}[leftmargin=2em, itemsep=0.1em]
        \item \textbf{Centro:} $C(-4,4)$.
        \item \textbf{Vértices:} $V(0,4)$ y $V'(-8,4)$; eje transverso $2a=8$.
        \item \textbf{Eje conjugado:} $2b=2$, sobre la recta $x=-4$.
        \item \textbf{Focos:} $F(-4+\sqrt{17},\,4)\approx(0{,}12;\,4)$ y
        $F'(-4-\sqrt{17},\,4)\approx(-8{,}12;\,4)$; distancia focal
        $2c=2\sqrt{17}\approx 8{,}25$.
        \item \textbf{Asíntotas:} $y-4=\pm\dfrac14(x+4)$, es decir
        \[
        y=\frac{x}{4}+5
        \qquad\text{y}\qquad
        y=-\frac{x}{4}+3 .
        \]
        \item \textbf{Excentricidad:}
        $e=\dfrac{\sqrt{17}}{4}\approx 1{,}03$.
        \item \textbf{Lado recto:} $\dfrac{2b^{2}}{a}=\dfrac12$.
    \end{itemize}

    \begin{center}
    \begin{tikzpicture}
        \begin{axis}[
            width=13cm, axis equal image,
            axis lines=middle,
            xlabel={$x$}, ylabel={$y$},
            xmin=-15, xmax=7.5, ymin=-0.5, ymax=8.5,
            xtick={-14,-12,-10,-8,-6,-4,-2,2,4,6},
            ytick={2,4,6,8},
            axis line style={-Latex},
            every axis x label/.style={at={(ticklabel* cs:1.0)}, anchor=west},
            every axis y label/.style={at={(ticklabel* cs:1.0)}, anchor=south},
        ]
            \addplot[dashed, OliveGreen, domain=-15:7.5] {x/4+5};
            \addplot[dashed, OliveGreen, domain=-15:7.5] {-x/4+3};
            \addplot[thick, blue, domain=-1.8:1.8, samples=120]
                ({-4+4*cosh(x)}, {4+sinh(x)});
            \addplot[thick, blue, domain=-1.8:1.8, samples=120]
                ({-4-4*cosh(x)}, {4+sinh(x)});
            \addplot[only marks, mark=*, mark size=1.6pt, black]
                coordinates {(-4,4) (0,4) (-8,4)};
            \addplot[only marks, mark=*, mark size=1.6pt, red]
                coordinates {(0.123,4) (-8.123,4)};
            \node[below, inner sep=5pt] at (axis cs:-4,4) {$C$};
            \node[above right, inner sep=1pt] at (axis cs:0,4) {$V$};
            \node[above left, inner sep=1pt]  at (axis cs:-8,4) {$V'$};
            \node[below right, red, inner sep=1pt] at (axis cs:0.123,4) {$F$};
            \node[below left, red, inner sep=1pt] at (axis cs:-8.123,4) {$F'$};
        \end{axis}
    \end{tikzpicture}
    \end{center}

\end{enumerate}

\volverindice
\newpage

% ======================== EJERCICIO 7 ========================

\section{Ecuación de la hipérbola a partir de sus condiciones}

Deduzca la ecuación de la hipérbola que cumple con las condiciones dadas y
grafíquela.

\begin{enumerate}[label=\alph*.-]

    \item Tiene focos en los puntos $(-3,0)$ y $(3,0)$, y la distancia entre sus
    vértices es $4$.

    \solucion

    \textbf{1. Centro y orientación.} El centro es el punto medio entre los
    focos:
    \[
    C=\left(\frac{-3+3}{2},\;\frac{0+0}{2}\right)=(0,0).
    \]
    Los focos están sobre el eje $x$, de modo que el \textbf{eje transverso es
    horizontal} y la ecuación tiene la forma
    \[
    \frac{x^{2}}{a^{2}}-\frac{y^{2}}{b^{2}}=1 .
    \]

    \textbf{2. Parámetros.} La semidistancia focal es $c=3$. La distancia entre
    los vértices es el eje transverso:
    \[
    2a=4\;\Longrightarrow\;a=2 .
    \]
    De la relación $c^{2}=a^{2}+b^{2}$:
    \[
    b^{2}=c^{2}-a^{2}=9-4=5
    \;\Longrightarrow\;
    b=\sqrt5\approx 2{,}24 .
    \]

    \textbf{3. Ecuación.}

    \[
    \boxed{
    \frac{x^{2}}{4}-\frac{y^{2}}{5}=1
    }
    \]

    \textbf{Elementos:} $C(0,0)$; vértices $V(2,0)$ y $V'(-2,0)$; focos
    $F(3,0)$ y $F'(-3,0)$; eje transverso $2a=4$; eje conjugado
    $2b=2\sqrt5\approx 4{,}47$; asíntotas
    $y=\pm\dfrac{\sqrt5}{2}\,x$; excentricidad
    $e=\dfrac{c}{a}=\dfrac32=1{,}5$.

    \begin{center}
    \begin{tikzpicture}
        \begin{axis}[
            width=9cm, axis equal image,
            axis lines=middle,
            xlabel={$x$}, ylabel={$y$},
            xmin=-7.5, xmax=7.5, ymin=-7, ymax=7,
            xtick={-6,-4,-2,2,4,6}, ytick={-6,-4,-2,2,4,6},
            axis line style={-Latex},
            every axis x label/.style={at={(ticklabel* cs:1.0)}, anchor=west},
            every axis y label/.style={at={(ticklabel* cs:1.0)}, anchor=south},
        ]
            \addplot[dashed, OliveGreen, domain=-7.5:7.5] {1.118034*x};
            \addplot[dashed, OliveGreen, domain=-7.5:7.5] {-1.118034*x};
            \addplot[thick, blue, domain=-1.8:1.8, samples=120]
                ({2*cosh(x)}, {2.236068*sinh(x)});
            \addplot[thick, blue, domain=-1.8:1.8, samples=120]
                ({-2*cosh(x)}, {2.236068*sinh(x)});
            \addplot[only marks, mark=*, mark size=1.6pt, black]
                coordinates {(0,0) (2,0) (-2,0)};
            \addplot[only marks, mark=*, mark size=1.6pt, red]
                coordinates {(3,0) (-3,0)};
            \node[below left, inner sep=3pt]  at (axis cs:0,0) {$C$};
            \node[above right, inner sep=1pt] at (axis cs:2,0) {$V$};
            \node[above left, inner sep=1pt]  at (axis cs:-2,0) {$V'$};
            \node[below right, red, inner sep=1pt] at (axis cs:3,0) {$F$};
            \node[below left, red, inner sep=1pt]  at (axis cs:-3,0) {$F'$};
        \end{axis}
    \end{tikzpicture}
    \end{center}

    \item Tiene centro en el punto $(3,-5)$, un vértice en el punto $(7,-5)$ y
    un foco en el punto $(8,-5)$.

    \solucion

    \textbf{1. Orientación.} El centro, el vértice y el foco tienen la misma
    ordenada $y=-5$, de modo que el \textbf{eje transverso es horizontal} y la
    ecuación tiene la forma
    \[
    \frac{(x-3)^{2}}{a^{2}}-\frac{(y+5)^{2}}{b^{2}}=1 .
    \]

    \textbf{2. Parámetros.} Medimos las distancias desde el centro:
    \[
    a=|7-3|=4,
    \qquad
    c=|8-3|=5 ,
    \]
    \[
    b^{2}=c^{2}-a^{2}=25-16=9
    \;\Longrightarrow\;
    b=3 .
    \]

    \textbf{3. Ecuación.}

    \[
    \boxed{
    \frac{(x-3)^{2}}{16}-\frac{(y+5)^{2}}{9}=1
    }
    \]

    \textbf{Elementos:} $C(3,-5)$; vértices $V(7,-5)$ y $V'(-1,-5)$; focos
    $F(8,-5)$ y $F'(-2,-5)$; eje transverso $2a=8$; eje conjugado $2b=6$;
    asíntotas
    \[
    y+5=\pm\frac34\,(x-3)
    \quad\Longrightarrow\quad
    y=\frac34x-\frac{29}{4}
    \;\;\text{y}\;\;
    y=-\frac34x-\frac{11}{4};
    \]
    excentricidad $e=\dfrac54=1{,}25$.

    \begin{center}
    \begin{tikzpicture}
        \begin{axis}[
            width=12cm, axis equal image,
            axis lines=middle,
            xlabel={$x$}, ylabel={$y$},
            xmin=-8.5, xmax=14.5, ymin=-12.5, ymax=2.5,
            xtick={-8,-6,-4,-2,2,4,6,8,10,12,14},
            ytick={-12,-10,-8,-6,-4,-2,2},
            axis line style={-Latex},
            every axis x label/.style={at={(ticklabel* cs:1.0)}, anchor=west},
            every axis y label/.style={at={(ticklabel* cs:1.0)}, anchor=south},
        ]
            \addplot[dashed, OliveGreen, domain=-8.5:14.5] {0.75*x-7.25};
            \addplot[dashed, OliveGreen, domain=-8.5:14.5] {-0.75*x-2.75};
            \addplot[thick, blue, domain=-1.5:1.5, samples=120]
                ({3+4*cosh(x)}, {-5+3*sinh(x)});
            \addplot[thick, blue, domain=-1.5:1.5, samples=120]
                ({3-4*cosh(x)}, {-5+3*sinh(x)});
            \addplot[only marks, mark=*, mark size=1.6pt, black]
                coordinates {(3,-5) (7,-5) (-1,-5)};
            \addplot[only marks, mark=*, mark size=1.6pt, red]
                coordinates {(8,-5) (-2,-5)};
            \node[below, inner sep=5pt] at (axis cs:3,-5) {$C$};
            \node[above left, inner sep=1pt]  at (axis cs:7,-5) {$V$};
            \node[above right, inner sep=1pt] at (axis cs:-1,-5) {$V'$};
            \node[below right, red, inner sep=1pt] at (axis cs:8,-5) {$F$};
            \node[below left, red, inner sep=1pt]  at (axis cs:-2,-5) {$F'$};
        \end{axis}
    \end{tikzpicture}
    \end{center}

    \item Tiene focos en los puntos $(-2,-1)$ y $(-2,-5)$ y longitud del eje
    imaginario $2$.

    \solucion

    \textbf{1. Centro y orientación.} El centro es el punto medio entre los
    focos:
    \[
    C=\left(\frac{-2-2}{2},\;\frac{-1+(-5)}{2}\right)=(-2,-3).
    \]
    Los focos tienen la misma abscisa, de modo que el \textbf{eje transverso es
    vertical} y la ecuación tiene la forma
    \[
    \frac{(y+3)^{2}}{a^{2}}-\frac{(x+2)^{2}}{b^{2}}=1 .
    \]

    \textbf{2. Parámetros.} La semidistancia focal es la distancia del centro a
    un foco:
    \[
    c=|-1-(-3)|=2 .
    \]
    El eje imaginario es el eje conjugado, de longitud $2b$:
    \[
    2b=2\;\Longrightarrow\;b=1 .
    \]
    De $c^{2}=a^{2}+b^{2}$:
    \[
    a^{2}=c^{2}-b^{2}=4-1=3
    \;\Longrightarrow\;
    a=\sqrt3\approx 1{,}73 .
    \]

    \textbf{3. Ecuación.}

    \[
    \boxed{
    \frac{(y+3)^{2}}{3}-\frac{(x+2)^{2}}{1}=1
    }
    \]

    \textbf{Elementos:} $C(-2,-3)$; vértices
    $V(-2,\,-3+\sqrt3)\approx(-2;\,-1{,}27)$ y
    $V'(-2,\,-3-\sqrt3)\approx(-2;\,-4{,}73)$; focos $F(-2,-1)$ y $F'(-2,-5)$;
    eje transverso $2a=2\sqrt3\approx 3{,}46$; eje conjugado $2b=2$;
    asíntotas
    \[
    y+3=\pm\frac{a}{b}\,(x+2)=\pm\sqrt3\,(x+2);
    \]
    excentricidad $e=\dfrac{c}{a}=\dfrac{2}{\sqrt3}
    =\dfrac{2\sqrt3}{3}\approx 1{,}15$.

    \begin{center}
    \begin{tikzpicture}
        \begin{axis}[
            width=8.5cm, height=13cm, axis equal image,
            axis lines=middle,
            xlabel={$x$}, ylabel={$y$},
            xmin=-5.2, xmax=1.2, ymin=-7.5, ymax=1.5,
            xtick={-5,-4,-3,-2,-1,1}, ytick={-7,-6,-5,-4,-3,-2,-1,1},
            axis line style={-Latex},
            every axis x label/.style={at={(ticklabel* cs:1.0)}, anchor=west},
            every axis y label/.style={at={(ticklabel* cs:1.0)}, anchor=south},
        ]
            \addplot[dashed, OliveGreen, domain=-5.2:1.2] {1.732051*(x+2)-3};
            \addplot[dashed, OliveGreen, domain=-5.2:1.2] {-1.732051*(x+2)-3};
            \addplot[thick, blue, domain=-1.3:1.3, samples=120]
                ({-2+sinh(x)}, {-3+1.732051*cosh(x)});
            \addplot[thick, blue, domain=-1.3:1.3, samples=120]
                ({-2+sinh(x)}, {-3-1.732051*cosh(x)});
            \addplot[only marks, mark=*, mark size=1.6pt, black]
                coordinates {(-2,-3) (-2,-1.268) (-2,-4.732)};
            \addplot[only marks, mark=*, mark size=1.6pt, red]
                coordinates {(-2,-1) (-2,-5)};
            \node[right, inner sep=5pt] at (axis cs:-2,-3) {$C$};
            \node[below left, inner sep=2pt]  at (axis cs:-2,-1.268) {$V$};
            \node[above left, inner sep=2pt]  at (axis cs:-2,-4.732) {$V'$};
            \node[above right, red, inner sep=2pt] at (axis cs:-2,-1) {$F$};
            \node[below right, red, inner sep=2pt] at (axis cs:-2,-5) {$F'$};
        \end{axis}
    \end{tikzpicture}
    \end{center}

\end{enumerate}

\volverindice
\newpage

% ======================== EJERCICIO 8 ========================

\section{Identificación de superficies cuádricas}

Dadas las siguientes ecuaciones identifique de qué superficie se trata, luego
mediante GeoGebra realice su gráfica, comprobando sus resultados.

\vspace{0.5em}
En cada caso llevamos la ecuación a su forma reducida y la comparamos con las
formas canónicas de la unidad. El criterio práctico es contar los signos de los
términos cuadráticos y observar si la tercera variable aparece elevada al
cuadrado o en primer grado:

\begin{itemize}[leftmargin=2em, itemsep=0.1em]
    \item tres cuadrados positivos $=1$: \emph{elipsoide};
    \item dos positivos y uno negativo $=1$: \emph{hiperboloide de una hoja};
    \item uno positivo y dos negativos $=1$: \emph{hiperboloide de dos hojas};
    \item dos cuadrados del mismo signo y la tercera variable de primer grado:
    \emph{paraboloide elíptico};
    \item dos cuadrados de signos opuestos y la tercera de primer grado:
    \emph{paraboloide hiperbólico};
    \item tres cuadrados $=0$ con signos mezclados: \emph{cono};
    \item casos degenerados: planos, rectas o un punto.
\end{itemize}

Las gráficas que siguen se generaron con \texttt{pgfplots} y reproducen las
superficies que se obtienen al ingresar cada ecuación en GeoGebra.

\begin{enumerate}[label=\alph*.-]

    \item $36x^{2}+9y^{2}-4z^{2}=36$

    \solucion

    Dividimos todo por $36$:
    \[
    \frac{36x^{2}}{36}+\frac{9y^{2}}{36}-\frac{4z^{2}}{36}=1
    \quad\Longrightarrow\quad
    \boxed{
    \frac{x^{2}}{1}+\frac{y^{2}}{4}-\frac{z^{2}}{9}=1
    }
    \]

    Hay \textbf{dos términos positivos y uno negativo}, igualados a $1$: se
    trata de un \textbf{hiperboloide de una hoja}, con centro en el origen y
    eje de simetría el \textbf{eje $z$} (el de la variable con signo negativo).
    Los semiejes son $1$ en $x$, $2$ en $y$ y $3$ en $z$.

    \textbf{Trazas:} con $z=k$ se obtiene
    $x^{2}+\dfrac{y^{2}}{4}=1+\dfrac{k^{2}}{9}$, una familia de elipses que
    crecen al alejarse del plano $xy$; con $x=0$ y con $y=0$ se obtienen
    hipérbolas.

    \begin{center}
    \begin{tikzpicture}
        \begin{axis}[
            width=8.5cm, view={120}{22},
            xlabel={$x$}, ylabel={$y$}, zlabel={$z$},
            colormap/viridis,
        ]
            \addplot3[surf, opacity=0.85, faceted color=black!25,
                      domain=0:360, y domain=-1.2:1.2,
                      samples=40, samples y=18]
                ({cosh(y)*cos(x)}, {2*cosh(y)*sin(x)}, {3*sinh(y)});
        \end{axis}
    \end{tikzpicture}
    \end{center}

    \item $x^{2}+z^{2}=9y$

    \solucion

    Despejamos la variable que aparece en primer grado:
    \[
    \boxed{
    y=\frac{x^{2}+z^{2}}{9}
    }
    \]

    Hay \textbf{dos cuadrados del mismo signo} ($x^{2}$ y $z^{2}$) y la tercera
    variable, $y$, en \textbf{primer grado}: es un \textbf{paraboloide
    elíptico}. Como los coeficientes de $x^{2}$ y $z^{2}$ son iguales, las
    secciones son circunferencias, de modo que se trata de un
    \textbf{paraboloide circular (de revolución)} de eje $y$, con vértice en el
    origen, que \textbf{abre hacia los $y$ positivos}.

    \textbf{Trazas:} con $y=k>0$ se obtiene $x^{2}+z^{2}=9k$, circunferencias
    de radio $3\sqrt{k}$; con $x=0$ resulta la parábola $z^{2}=9y$ y con $z=0$
    la parábola $x^{2}=9y$.

    \begin{center}
    \begin{tikzpicture}
        \begin{axis}[
            width=8.5cm, view={120}{22},
            xlabel={$x$}, ylabel={$y$}, zlabel={$z$},
            colormap/viridis,
        ]
            \addplot3[surf, opacity=0.85, faceted color=black!25,
                      domain=0:4, y domain=0:360,
                      samples=18, samples y=40]
                ({x*cos(y)}, {x^2/9}, {x*sin(y)});
        \end{axis}
    \end{tikzpicture}
    \end{center}

    \item $9x^{2}+y^{2}+z^{2}=36$

    \solucion

    Dividimos por $36$:
    \[
    \boxed{
    \frac{x^{2}}{4}+\frac{y^{2}}{36}+\frac{z^{2}}{36}=1
    }
    \]

    Los \textbf{tres términos son cuadrados positivos} igualados a $1$: es un
    \textbf{elipsoide} con centro en el origen y semiejes
    \[
    a=2\;(\text{en }x),
    \qquad
    b=6\;(\text{en }y),
    \qquad
    c=6\;(\text{en }z).
    \]

    Como dos de los semiejes coinciden ($b=c=6$), es además un
    \textbf{elipsoide de revolución} alrededor del eje $x$.

    \textbf{Trazas:} con $x=k$ ($|k|<2$) se obtienen circunferencias; con $y=k$
    o $z=k$, elipses.

    \begin{center}
    \begin{tikzpicture}
        \begin{axis}[
            width=8.5cm, view={120}{22},
            xlabel={$x$}, ylabel={$y$}, zlabel={$z$},
            colormap/viridis,
        ]
            \addplot3[surf, opacity=0.85, faceted color=black!25,
                      domain=0:180, y domain=0:360,
                      samples=24, samples y=40]
                ({2*sin(x)*cos(y)}, {6*sin(x)*sin(y)}, {6*cos(x)});
        \end{axis}
    \end{tikzpicture}
    \end{center}

    \item $y^{2}-9x^{2}=9z$

    \solucion

    Despejamos la variable de primer grado:
    \[
    \boxed{
    z=\frac{y^{2}}{9}-x^{2}
    }
    \]

    Hay \textbf{dos cuadrados de signos opuestos} y la tercera variable, $z$, en
    \textbf{primer grado}: es un \textbf{paraboloide hiperbólico}, la
    característica superficie en forma de \emph{silla de montar}, con punto
    silla en el origen.

    \textbf{Trazas:} con $z=0$ resulta $y^{2}=9x^{2}$, es decir el par de
    rectas $y=\pm 3x$; con $x=0$ la parábola $z=\dfrac{y^{2}}{9}$, que abre
    hacia arriba; y con $y=0$ la parábola $z=-x^{2}$, que abre hacia abajo.

    \begin{center}
    \begin{tikzpicture}
        \begin{axis}[
            width=8.5cm, view={130}{22},
            xlabel={$x$}, ylabel={$y$}, zlabel={$z$},
            colormap/viridis,
        ]
            \addplot3[surf, opacity=0.85, faceted color=black!25,
                      domain=-3:3, y domain=-6:6, samples=28]
                {y^2/9 - x^2};
        \end{axis}
    \end{tikzpicture}
    \end{center}

    \item $x^{2}-z^{2}=0$

    \solucion

    Se trata de una \textbf{cuádrica degenerada}. Factorizamos como diferencia
    de cuadrados:
    \[
    x^{2}-z^{2}=0
    \quad\Longrightarrow\quad
    (x-z)(x+z)=0
    \quad\Longrightarrow\quad
    \boxed{
    x-z=0
    \;\;\text{ó}\;\;
    x+z=0
    }
    \]

    La ecuación no impone ninguna condición sobre $y$, de modo que cada factor
    describe un \textbf{plano} que contiene al eje $y$. La superficie es
    entonces el \textbf{par de planos} $z=x$ y $z=-x$, que
    \textbf{se cortan en el eje $y$} formando ángulos de $45^{\circ}$ con los
    planos coordenados.

    \begin{center}
    \begin{tikzpicture}
        \begin{axis}[
            width=8.5cm, view={120}{22},
            xlabel={$x$}, ylabel={$y$}, zlabel={$z$},
        ]
            \addplot3[surf, opacity=0.7, faceted color=black!30,
                      color=blue!45,
                      domain=-3:3, y domain=-3:3, samples=12]
                {x};
            \addplot3[surf, opacity=0.7, faceted color=black!30,
                      color=red!45,
                      domain=-3:3, y domain=-3:3, samples=12]
                {-x};
        \end{axis}
    \end{tikzpicture}
    \end{center}

    \item $x^{2}+3y^{2}+7z^{2}=0$

    \solucion

    Los tres coeficientes son \textbf{positivos} y la suma está igualada a
    \textbf{cero}. Como cada sumando cumple
    \[
    x^{2}\ge 0,
    \qquad
    3y^{2}\ge 0,
    \qquad
    7z^{2}\ge 0,
    \]
    la suma sólo puede anularse si los tres sumandos son simultáneamente nulos:
    \[
    x^{2}=0,\quad 3y^{2}=0,\quad 7z^{2}=0
    \quad\Longrightarrow\quad
    x=y=z=0 .
    \]

    \[
    \boxed{
    \text{La ecuación se reduce a un único punto: el origen }O(0,0,0).
    }
    \]

    Es una \textbf{cuádrica degenerada}: puede pensarse como el caso límite de
    un elipsoide cuyos tres semiejes tienden a cero (en algunos textos se la
    llama \emph{elipsoide imaginario} o \emph{punto}).

    \begin{center}
    \begin{tikzpicture}
        \begin{axis}[
            width=8cm, view={120}{22},
            xlabel={$x$}, ylabel={$y$}, zlabel={$z$},
            xmin=-2, xmax=2, ymin=-2, ymax=2, zmin=-2, zmax=2,
        ]
            \addplot3[only marks, mark=*, mark size=2.5pt, red]
                coordinates {(0,0,0)};
            \node[above right, red] at (axis cs:0,0,0) {$O(0,0,0)$};
        \end{axis}
    \end{tikzpicture}
    \end{center}

    \item $x^{2}-y^{2}+z^{2}=0$

    \solucion

    Pasamos el término negativo al otro miembro:
    \[
    \boxed{
    y^{2}=x^{2}+z^{2}
    }
    \]

    Hay \textbf{tres términos cuadráticos con signos mezclados igualados a
    cero}: se trata de un \textbf{cono}. Como los coeficientes de $x^{2}$ y
    $z^{2}$ son iguales, es un \textbf{cono circular recto}, con
    \textbf{vértice en el origen} y \textbf{eje de simetría el eje $y$}.
    Equivale a las dos hojas
    \[
    y=\sqrt{x^{2}+z^{2}}
    \qquad\text{y}\qquad
    y=-\sqrt{x^{2}+z^{2}} .
    \]

    \textbf{Trazas:} con $y=k$ se obtiene $x^{2}+z^{2}=k^{2}$, circunferencias
    de radio $|k|$ que crecen linealmente al alejarse del vértice; con $x=0$
    resulta $y=\pm z$ y con $z=0$ resulta $y=\pm x$, en ambos casos un par de
    rectas que se cortan en el origen.

    \begin{center}
    \begin{tikzpicture}
        \begin{axis}[
            width=8.5cm, view={120}{22},
            xlabel={$x$}, ylabel={$y$}, zlabel={$z$},
            colormap/viridis,
        ]
            \addplot3[surf, opacity=0.85, faceted color=black!25,
                      domain=0:3, y domain=0:360,
                      samples=14, samples y=40]
                ({x*cos(y)}, {x}, {x*sin(y)});
            \addplot3[surf, opacity=0.85, faceted color=black!25,
                      domain=0:3, y domain=0:360,
                      samples=14, samples y=40]
                ({x*cos(y)}, {-x}, {x*sin(y)});
        \end{axis}
    \end{tikzpicture}
    \end{center}

\end{enumerate}

\volverindice
\newpage

\end{document}
